\section{The input keys that select physics}
\label{sec:keys}

Runtime configuration is by the presence of a key, not by a compile flag
(Sec.~\ref{sec:intro:what}).  Three companion files are read at startup and
their mere presence is itself a switch:
\texttt{metals.inp} turns the trace metals on, \texttt{opacity.inp} the
tabulated continuum opacities, and \texttt{base.inp} the lower-atmosphere
handoff.

Table~\ref{tab:keys} lists the keys that select physics or the numerical
scheme, with the default the code takes when the key is absent.  The defaults
were read from \texttt{parameters.f90} and the parsing from
\texttt{input\_read.f90}; keys that only name a file, a path or an output
option are omitted.  Four lines of the core block are mandatory and have no
default: they are marked as such.

{\small
\begin{longtable}{@{}>{\raggedright\arraybackslash}p{0.245\textwidth}>{\raggedright\arraybackslash}p{0.155\textwidth}>{\raggedright\arraybackslash}p{0.45\textwidth}@{}}
\caption{Physics and algorithm keys of \texttt{input.inp}, with the value the
code takes when the key is absent.}
\label{tab:keys}\\
\toprule
Key & Default & What it selects, and where \\
\midrule
\endfirsthead
\multicolumn{3}{l}{\small\itshape Table \thetable{} continued}\\
\toprule
Key & Default & What it selects, and where \\
\midrule
\endhead
\bottomrule
\endfoot

\multicolumn{3}{l}{\itshape Geometry, grid and the one-dimensional approximation}\\
\midrule
\texttt{2D approximate method:} & mandatory &
  \texttt{Mdot/4}, \texttt{Rate/2}, \texttt{Rate/4} or \texttt{alpha}: the
  dayside dilution of every stellar beam and the factor on $\dot M$
  (Table~\ref{tab:2d}; \texttt{parameters.f90}, \texttt{utilities.f90}). \\
\texttt{Grid type:} & mandatory &
  \texttt{Uniform}, \texttt{Mixed} or \texttt{Stretched}
  (\texttt{define\_grid.f90}). \\
\texttt{Grid cells:} & 500 &
  Number of physical cells; ghosts are added on top
  (\texttt{define\_grid.f90}). \\
\texttt{Base grid [dr,cells]:} & $2\times10^{-4}$, 50 &
  Width and count of the uniform base region of the \texttt{Mixed} grid
  (\texttt{define\_grid.f90}). \\
\texttt{Domain mode:} & Roche &
  \texttt{Spherical} drops the tidal and centrifugal terms of
  \eqref{eq:roche} (\texttt{grav\_field.f90}). \\
\texttt{Outer radius [R\_p]:} & L1 radius &
  Outer boundary; required in spherical mode
  (\texttt{input\_read.f90}). \\
\midrule
\multicolumn{3}{l}{\itshape Discretization}\\
\midrule
\texttt{Numerical flux:} & mandatory &
  \texttt{HLLC}, \texttt{ROE} or \texttt{LLF}
  (\texttt{Num\_Fluxes.f90}); \texttt{ROE} is refused with the molecular
  ladder equation of state. \\
\texttt{Reconstruction scheme:} & mandatory &
  \texttt{PLM}, \texttt{WENO3} or the two-stage \texttt{PLM+WENO3}
  (\texttt{Reconstruction.f90}). \\
\texttt{Reconstruction continuation:} & off &
  $\Delta\lambda$ of the homotopy \eqref{eq:homotopy} between the two
  operators (\texttt{steady\_residual.f90}). \\
\texttt{CFL:} & 0.6 &
  CFL number of \eqref{eq:cfl} (\texttt{eval\_dt.f90}). \\
\texttt{Time stepping:} & global &
  \texttt{Local} gives every cell its own pseudo-time step
  (\texttt{eval\_dt.f90}). \\
\texttt{Well balanced:} & off &
  Carries the departure from the cell's own hydrostatic equilibrium instead of
  the state (\texttt{Reconstruction.f90}, \texttt{RK\_rhs.f90}). \\
\texttt{Low-Mach damping:} & off, $M_{\rm th}=10^{-3}$ &
  Gated fourth-difference dissipation \eqref{eq:lowmach} of the contact mode
  (\texttt{low\_mach\_dissipation.f90}). \\
\texttt{Shapiro filter:} & off, every 4 steps &
  1-2-1 low-pass filter on the marched state
  (\texttt{Apply\_BC.f90}). \\
\texttt{Caloric EOS:} & \texttt{ladder} &
  \texttt{monatomic} forces $\gamma = 5/3$ everywhere
  (\texttt{caloric\_eos.f90}). \\
\texttt{Energy solver:} & semi-implicit &
  \texttt{Explicit} reverts the source update to forward Euler
  (\texttt{energy\_semi\_implicit.f90}). \\
\midrule
\multicolumn{3}{l}{\itshape Boundary and initial conditions}\\
\midrule
\texttt{Base BC:} & \texttt{density} &
  \texttt{pressure <microbar>} sets the base level from a pressure instead of
  a density (\texttt{input\_read.f90}, \texttt{base\_boundary.f90}). \\
\texttt{IC mode:} & \texttt{cold} &
  \texttt{cold}, \texttt{transonic}, \texttt{hot\_parker}, \texttt{auto} or
  \texttt{windae} (\texttt{set\_IC.f90}). \\
\texttt{Transonic IC:} & off &
  Older synonym of \texttt{IC mode: transonic}. \\
\texttt{Hot Parker IC:} & off, $T_{\rm wind}=10^4$~K &
  Older synonym of \texttt{IC mode: hot\_parker}. \\
\texttt{Restart intent:} & relaxation &
  What a loaded state is continued as: \texttt{trajectory},
  \texttt{relaxation} or \texttt{stationary} (\texttt{load\_IC.f90}). \\
\midrule
\multicolumn{3}{l}{\itshape Convergence and the stationary solvers}\\
\midrule
\texttt{du\_th [PLM,WENO3]:} & $10^{-3}$ &
  Threshold on the mass-flux spread \eqref{eq:du}; a second value makes the
  run two-stage (\texttt{EXHALE\_main.f90}). \\
\texttt{Escape radius [R\_p]:} & mandatory &
  Inner edge of the window \eqref{eq:du} is taken over
  (\texttt{define\_grid.f90}). \\
\texttt{Resid tol:} & off &
  Residual convergence instead of the flux spread; the solvers use $10^{-5}$
  when it is absent (\texttt{steady\_residual.f90}). \\
\texttt{Flux spread tol:} & $2\times10^{-5}$, $r_{\rm flux}=1.2$ &
  The flux gate of the acceptance test
  (\texttt{steady\_residual.f90}). \\
\texttt{Solver:} & marching only &
  \texttt{Newton} hands over to the stationary solve once the spread falls
  below $10^{-2}$ (\texttt{steady\_newton.f90}). \\
\texttt{Stall:} & $10^{-6}$, 2000 steps &
  Plateau detector of the marching stop (\texttt{EXHALE\_main.f90}). \\
\texttt{Level tol:} & off &
  Mass-flux level stability required beside every marching stop. \\
\texttt{Max steps:} & $10^{6}$ &
  Hard cap on marching iterations. \\
\texttt{Newton solver:} & on &
  \texttt{False} reverts the composition solve to MINPACK \texttt{hybrd1}
  (\texttt{newton\_solver.f90}). \\
\texttt{Brent solver:} & on &
  \texttt{False} reverts the post-process temperature solve
  (\texttt{T\_equation.f90}). \\
\midrule
\multicolumn{3}{l}{\itshape Atomic hydrogen and helium}\\
\midrule
\texttt{Include He23S?} & on &
  The $\mathrm{He}\,2^3S$ metastable as a solved level
  (\texttt{System\_HeH\_TR*.f90}). \\
\texttt{ATES\_photoionization\_rate} & off &
  Reverts He\,I photoionization to the legacy ATES fit
  (\texttt{cross\_sec.f90}). \\
\texttt{Legacy\_HHe\_rates} & off &
  Reverts H/He recombination and collisional ionization to the ATES fits
  (\texttt{util\_ion\_eq.f90}). \\
\texttt{Atomic rate set:} & \texttt{default} &
  \texttt{Koskinen2022} swaps four atomic H/He rate coefficients
  (\texttt{util\_ion\_eq.f90}). \\
\texttt{Secondary\_ionization} & on, staged &
  Secondary ionization by fast photoelectrons; \texttt{Immediate} applies it
  from step 0 (\texttt{electron\_energy\_degradation.f90}). \\
\texttt{He\_rec\_coupling} & on &
  He\,II recombination radiation ionizing H\,I
  (\texttt{util\_ion\_eq.f90}). \\
\texttt{He\_H\_charge\_exchange} & on &
  The He/H charge-exchange pair (\texttt{charge\_exchange.f90}). \\
\texttt{Photoelectron heating:} & \texttt{excess} &
  Whether a photoionization deposits $h\nu - E_{\rm th}$ or a stated fraction
  of $h\nu$ (\texttt{util\_ion\_eq.f90}). \\
\texttt{Deexc heat} & on &
  Collisional de-excitation heating of $\mathrm{H}(n{=}2)$
  (\texttt{excited\_hydrogen.f90}). \\
\midrule
\multicolumn{3}{l}{\itshape Radiation field}\\
\midrule
\texttt{Spectrum type:} & mandatory &
  \texttt{Load}, \texttt{Power-law}, \texttt{Planck} or
  \texttt{Monochromatic}: the spectrum of the whole photon grid
  (\texttt{sed\_read.f90}, \texttt{J\_inc.f90}). \\
\texttt{Stellar Teff:} + \texttt{Stellar radius:} & 0 (off) &
  Enable the diluted-blackbody Balmer continuum and the $\mathrm{H}(n{=}2)$
  model (\texttt{excited\_hydrogen.f90}). \\
\texttt{Stellar Lya flux:} & 0 &
  Stellar Ly-$\alpha$ at the orbit; also band B2 of the oxygen photolysis
  (\texttt{lya\_rt.f90}). \\
\texttt{Jlya escape-prob} & parameterized &
  The Neufeld escape-probability Ly-$\alpha$ transfer
  (\texttt{lya\_rt.f90}). \\
\texttt{Jlya RT file:} & none &
  Import a $J_{\mathrm{Ly}\alpha}(r)$ profile computed elsewhere. \\
\texttt{Lya stellar halfwidth:} & 70 km/s &
  Half-width of the stellar Ly-$\alpha$ line. \\
\texttt{Lya stellar boost:} & 5.0 &
  Bounded buildup of the trapped stellar beam. \\
\texttt{Lya absorbing bottom} & reflecting &
  Closes the Ly-$\alpha$ domain from below with an absorber. \\
\texttt{Stellar LW flux:} & 0 &
  Lyman-Werner band flux driving $\mathrm{H}_2$ photodissociation
  (\texttt{lyman\_werner.f90}). \\
\texttt{Stellar FUV B3/B4 flux:} & 0 &
  The two longer FUV bands of the oxygen photolysis
  (\texttt{water\_photolysis.f90}). \\
\midrule
\multicolumn{3}{l}{\itshape Molecules, metals and the lower atmosphere}\\
\midrule
\texttt{metals.inp} present & metals off &
  Trace metals solved in the same coupled system as H and He
  (\texttt{System\_HeH\_metals.f90}). \\
\quad \texttt{eos\_metals} & 1 &
  Metals in the bulk mass, electron and particle budget
  (\texttt{composition.f90}). \\
\quad \texttt{cno\_cool} & 1 (CHIANTI) &
  0 reverts C/N/O line cooling to the legacy analytic fits
  (\texttt{Cool\_coeff.f90}). \\
\texttt{Molecular chemistry:} & off &
  The $\mathrm{H_2}/\mathrm{H_2^+}/\mathrm{H_3^+}/\mathrm{HeH^+}$ network
  (\texttt{System\_HeH\_mol.f90}, \texttt{mol\_rates.f90}). \\
\texttt{Oxygen chemistry:} & off &
  OH, $\mathrm{H_2O}$ and CO with their photolysis
  (\texttt{oxygen\_rates.f90}, \texttt{water\_photolysis.f90}). \\
\texttt{Molecular carrier transport:} & follows oxygen chemistry &
  Vertical transport of the molecular carriers
  (\texttt{diffusive\_photochemistry.f90}). \\
\texttt{Coupled carrier solve:} & off &
  $n(\mathrm{H_2})$ as a fourth Newton unknown per cell instead of an
  alternation (\texttt{steady\_newton.f90}). \\
\texttt{Ionization transport:} & off &
  $\mathrm{H^+}$ carried with the flow as a fifth carrier. \\
\texttt{Molecular base:} & off &
  Removes the H nuclei bound into $\mathrm{H}_2$ from the base particle count
  \eqref{eq:ntotbc} (\texttt{composition.f90}). \\
\texttt{Molecular reaction heat:} & on &
  Deposits the energy the collisional molecular reactions release
  (\texttt{molecular\_reaction\_heat.f90}). \\
\texttt{Molecular IR bands:} & off &
  $\mathrm{H_2}$, $\mathrm{H_2O}$ and CO infrared line cooling
  (\texttt{molecular\_infrared\_cooling.f90}). \\
\texttt{Base IR field:} & off &
  The infrared coolants see the thermal radiation of the atmosphere below the
  base instead of vacuum. \\
\texttt{Lower atmosphere:} & none &
  \texttt{analytic} or \texttt{vulcan} lower-atmosphere pre-step
  (\texttt{lower\_column.f90}). \\
\texttt{Lower atmosphere profile:} & none &
  The lower atmosphere handed over as a table over pressure
  (\texttt{lower\_atmosphere\_profile.f90}). \\
\texttt{Lower column:} & off &
  1-bar radius of the analytic column. \\
\midrule
\multicolumn{3}{l}{\itshape Diffusion and molecular transport}\\
\midrule
\texttt{He\_diffusion} & off &
  He/H diffusive separation as a transported element
  (\texttt{binary\_element\_diffusion.f90}). \\
\texttt{He\_metal\_diffusion} & off &
  Extends the same operator to the trace metals. \\
\texttt{He\_Kzz} & 0 &
  Constant eddy diffusion coefficient where no profile states one. \\
\texttt{He\_alphaT} & 0 &
  Thermal-diffusion factor. \\
\texttt{He\_ambipolar} & on &
  Ambipolar term of the element diffusion. \\
\texttt{Viscosity:} & off &
  Radial viscous force and its dissipation \eqref{eq:viscous}
  (\texttt{viscous\_conduction.f90}). \\
\texttt{Conduction:} & off &
  Heat conduction with $\kappa(T)$ of \eqref{eq:transport_coeff}. \\
\end{longtable}
}

\begin{thebibliography}{99}

\bibitem[Abrahamsson et al.(2007)]{Abrahamsson2007}
Abrahamsson, E., Krems, R.~V., \& Dalgarno, A. 2007, ApJ, 654, 1171

\bibitem[Allgower \& Georg(2003)]{Allgower2003}
Allgower, E.~L., \& Georg, K. 2003, Introduction to Numerical Continuation
Methods, SIAM Classics in Applied Mathematics 45 (Philadelphia: SIAM)

\bibitem[Backx et al.(1976)]{Backx1976}
Backx, C., Wight, G.~R., \& Van der Wiel, M.~J. 1976, J. Phys. B, 9, 315

\bibitem[Badnell(2006)]{Badnell2006}
Badnell, N.~R. 2006, ApJS, 167, 334

\bibitem[Banks \& Kockarts(1973)]{Banks1973}
Banks, P.~M., \& Kockarts, G. 1973, Aeronomy (New York: Academic Press)

\bibitem[Barinovs et al.(2005)]{Barinovs2005}
Barinovs, \v{G}., van Hemert, M.~C., Krems, R., \& Dalgarno, A. 2005,
ApJ, 620, 537

\bibitem[Batten et al.(1997)]{Batten1997}
Batten, P., Clarke, N., Lambert, C., \& Causon, D.~M. 1997,
SIAM J. Sci. Comput., 18, 1553

\bibitem[Biassoni et al.(2024)]{Biassoni2024}
Biassoni, F., Caldiroli, A., Gallo, E., et al. 2024, A\&A, 682, A115

\bibitem[Caldiroli et al.(2021)]{Caldiroli2021}
Caldiroli, A., Haardt, F., Gallo, E., et al. 2021, A\&A, 655, A30

\bibitem[Cecchi-Pestellini et al.(2009)]{Cecchi2009}
Cecchi-Pestellini, C., Ciaravella, A., Micela, G., \& Penz, T. 2009,
A\&A, 496, 863

\bibitem[Christie et al.(2013)]{Christie2013}
Christie, D., Arras, P., \& Li, Z.-Y. 2013, ApJ, 772, 144

\bibitem[Chung et al.(1993)]{Chung1993}
Chung, Y.~M., Lee, E.-M., Masuoka, T., \& Samson, J.~A.~R. 1993,
J. Chem. Phys., 99, 885

\bibitem[Cohen \& Lane(1971)]{Cohen1971}
Cohen, J.~S., \& Lane, N.~F. 1971, Chem. Phys. Lett., 10, 623

\bibitem[Cohen \& Lane(1977)]{Cohen1977}
Cohen, J.~S., \& Lane, N.~F. 1977, J. Chem. Phys., 66, 586

\bibitem[Dalgarno et al.(1999)]{Dalgarno1999}
Dalgarno, A., Yan, M., \& Liu, W. 1999, ApJS, 125, 237

\bibitem[Draine(2011)]{Draine2011}
Draine, B.~T. 2011, Physics of the Interstellar and Intergalactic Medium
(Princeton: Princeton Univ. Press)

\bibitem[Draine \& Bertoldi(1996)]{Draine1996}
Draine, B.~T., \& Bertoldi, F. 1996, ApJ, 468, 269

\bibitem[Drake et al.(1969)]{Drake1969}
Drake, G.~W., Victor, G.~A., \& Dalgarno, A. 1969, Phys. Rev., 180, 25

\bibitem[Einfeldt et al.(1991)]{Einfeldt1991}
Einfeldt, B., Munz, C.-D., Roe, P.~L., \& Sj\"ogreen, B. 1991,
J. Comput. Phys., 92, 273

\bibitem[Garc\'ia Mu\~noz(2025)]{GarciaMunoz2025}
Garc\'ia Mu\~noz, A. 2025, A\&A, 698, A199

\bibitem[Gordon \& McBride(1994)]{Gordon1994}
Gordon, S., \& McBride, B.~J. 1994, Computer Program for Calculation of
Complex Chemical Equilibrium Compositions and Applications, NASA RP-1311

\bibitem[Gottlieb \& Shu(1998)]{Gottlieb1998}
Gottlieb, S., \& Shu, C.-W. 1998, Math. Comput., 67, 73

\bibitem[Harrington(1973)]{Harrington1973}
Harrington, J.~P. 1973, MNRAS, 162, 43

\bibitem[Hu et al.(2013)]{Hu2013}
Hu, X.~Y., Adams, N.~A., \& Shu, C.-W. 2013, J. Comput. Phys., 242, 169

\bibitem[Huang et al.(2023)]{Huang2023}
Huang, C., Koskinen, T., Lavvas, P., \& Fossati, L. 2023, ApJ, 951, 123

\bibitem[Hui \& Gnedin(1997)]{Hui1997}
Hui, L., \& Gnedin, N.~Y. 1997, MNRAS, 292, 27

\bibitem[Jameson et al.(1981)]{Jameson1981}
Jameson, A., Schmidt, W., \& Turkel, E. 1981, AIAA Paper 81-1259

\bibitem[K\"appeli \& Mishra(2014)]{Kaeppeli2014}
K\"appeli, R., \& Mishra, S. 2014, J. Comput. Phys., 259, 199

\bibitem[K\"appeli \& Mishra(2016)]{Kaeppeli2016}
K\"appeli, R., \& Mishra, S. 2016, A\&A, 587, A94

\bibitem[Knoll \& Keyes(2004)]{Knoll2004}
Knoll, D.~A., \& Keyes, D.~E. 2004, J. Comput. Phys., 193, 357

\bibitem[Koskinen et al.(2013)]{Koskinen2013}
Koskinen, T.~T., Harris, M.~J., Yelle, R.~V., \& Lavvas, P. 2013,
Icarus, 226, 1678

\bibitem[Koskinen et al.(2022)]{Koskinen2022}
Koskinen, T.~T., Lavvas, P., Huang, C., et al. 2022, ApJ, 929, 52

\bibitem[Krause(1979)]{Krause1979}
Krause, M.~O. 1979, J. Phys. Chem. Ref. Data, 8, 307

\bibitem[Kurganov \& Tadmor(2000)]{Kurganov2000}
Kurganov, A., \& Tadmor, E. 2000, J. Comput. Phys., 160, 241

\bibitem[Locci et al.(2022)]{Locci2022}
Locci, D., Cecchi-Pestellini, C., \& Micela, G. 2022, PSJ, 3, 1

\bibitem[Mao \& Kaastra(2016)]{Mao2016}
Mao, J., \& Kaastra, J. 2016, A\&A, 587, A84

\bibitem[Miller et al.(2013)]{Miller2013}
Miller, S., Stallard, T., Tennyson, J., \& Melin, H. 2013,
J. Phys. Chem. A, 117, 9770

\bibitem[Morgner \& Niehaus(1979)]{Morgner1979}
Morgner, H., \& Niehaus, A. 1979, J. Phys. B, 12, 1805

\bibitem[Neufeld(1990)]{Neufeld1990}
Neufeld, D.~A. 1990, ApJ, 350, 216

\bibitem[Norcross(1971)]{Norcross1971}
Norcross, D.~W. 1971, J. Phys. B, 4, 652

\bibitem[Odert et al.(2020)]{Odert2020}
Odert, P., Erkaev, N.~V., Lammer, H., et al. 2020, A\&A, 638, A49

\bibitem[Oka \& Epp(2004)]{Oka2004}
Oka, T., \& Epp, E. 2004, ApJ, 613, 349

\bibitem[Oklop\v{c}i\'c \& Hirata(2018)]{Oklopcic2018}
Oklop\v{c}i\'c, A., \& Hirata, C.~M. 2018, ApJ, 855, L11

\bibitem[Perthame \& Shu(1996)]{Perthame1996}
Perthame, B., \& Shu, C.-W. 1996, Numer. Math., 73, 119

\bibitem[Poinsot \& Lele(1992)]{Poinsot1992}
Poinsot, T.~J., \& Lele, S.~K. 1992, J. Comput. Phys., 101, 104

\bibitem[Robinson \& Catling(2012)]{Robinson2012}
Robinson, T.~D., \& Catling, D.~C. 2012, ApJ, 757, 104

\bibitem[Roe(1981)]{Roe1981}
Roe, P.~L. 1981, J. Comput. Phys., 43, 357

\bibitem[Roueff et al.(2019)]{Roueff2019}
Roueff, E., Abgrall, H., Czachorowski, P., et al. 2019, A\&A, 630, A58

\bibitem[Rusanov(1961)]{Rusanov1961}
Rusanov, V.~V. 1961, J. Comput. Math. Phys. USSR, 1, 267

\bibitem[Salz et al.(2016)]{Salz2016}
Salz, M., Czesla, S., Schneider, P.~C., \& Schmitt, J.~H.~M.~M. 2016,
A\&A, 586, A75

\bibitem[Samson \& Haddad(1994)]{Samson1994}
Samson, J.~A.~R., \& Haddad, G.~N. 1994, J. Opt. Soc. Am. B, 11, 277

\bibitem[Sekiya et al.(1980)]{Sekiya1980}
Sekiya, M., Nakazawa, K., \& Hayashi, C. 1980, Prog. Theor. Phys., 64, 1968

\bibitem[Shull \& van Steenberg(1985)]{Shull1985}
Shull, J.~M., \& van Steenberg, M.~E. 1985, ApJ, 298, 268

\bibitem[Taylor et al.(2025)]{Taylor2025}
Taylor, J., Oklop\v{c}i\'c, A., Rustamkulov, Z., et al. 2025, ApJ, 989, 68

\bibitem[Thompson(1987)]{Thompson1987}
Thompson, K.~W. 1987, J. Comput. Phys., 68, 1

\bibitem[Toro et al.(1994)]{Toro1994}
Toro, E.~F., Spruce, M., \& Speares, W. 1994, Shock Waves, 4, 25

\bibitem[Toro(2009)]{Toro2009}
Toro, E.~F. 2009, Riemann Solvers and Numerical Methods for Fluid Dynamics,
3rd edn. (Berlin: Springer)

\bibitem[van Hoof et al.(2014)]{vanHoof2014}
van Hoof, P.~A.~M., Williams, R.~J.~R., Volk, K., et al. 2014,
MNRAS, 444, 420

\bibitem[Verner \& Yakovlev(1995)]{Verner1995}
Verner, D.~A., \& Yakovlev, D.~G. 1995, A\&AS, 109, 125

\bibitem[Verner et al.(1996)]{Verner1996}
Verner, D.~A., Ferland, G.~J., Korista, K.~T., \& Yakovlev, D.~G. 1996,
ApJ, 465, 487

\bibitem[Visscher et al.(2006)]{Visscher2006}
Visscher, C., Lodders, K., \& Fegley, B., Jr. 2006, ApJ, 648, 1181

\bibitem[Visser et al.(2009)]{Visser2009}
Visser, R., van Dishoeck, E.~F., \& Black, J.~H. 2009, A\&A, 503, 323

\bibitem[Voronov(1997)]{Voronov1997}
Voronov, G.~S. 1997, At. Data Nucl. Data Tables, 65, 1

\bibitem[Watson et al.(1981)]{Watson1981}
Watson, A.~J., Donahue, T.~M., \& Walker, J.~C.~G. 1981, Icarus, 48, 150

\bibitem[Yamaleev \& Carpenter(2009)]{Yamaleev2009}
Yamaleev, N.~K., \& Carpenter, M.~H. 2009, J. Comput. Phys., 228, 3025

\bibitem[Yan et al.(1998)]{Yan1998}
Yan, M., Sadeghpour, H.~R., \& Dalgarno, A. 1998, ApJ, 496, 1044

\bibitem[Yan \& Babb(2023)]{YanBabb2023}
Yan, P.-G., \& Babb, J.~F. 2023, MNRAS, 518, 6004

\bibitem[Zhang \& Shu(2010a)]{Zhang2010}
Zhang, X., \& Shu, C.-W. 2010a, J. Comput. Phys., 229, 8918

\bibitem[Zhang \& Shu(2010b)]{Zhang2010b}
Zhang, X., \& Shu, C.-W. 2010b, J. Comput. Phys., 229, 3091


\end{thebibliography}
